\section{Feedback og stationær fejl}
\label{sec:feedback}

\subsection{Blokdiagramreduktion / feedback path, closed-loop transfer}

\paragraph{Brug når}
Et diagram viser flere forward branches, sensorfeedback eller en controller i
feedbackgrenen.

\paragraph{Input fra opgaven}
Alle blokke, summationsfortegn og det præcise input/outputpar der efterspørges.

\paragraph{Metode}
\begin{enumerate}
\item Skriv hvilket input/outputpar opgaven spørger om, fx \(Y/R\), \(E/R\) eller
      \(Y/D\). Alle andre uafhængige input sættes til nul; det betyder at reference,
      disturbance og noise behandles én ad gangen med superposition.
\item Sæt et symbol på hver vigtig ledning i diagrammet: \(e,u,y\) og eventuelle
      mellemled. For hvert summationspunkt skriver du en ligning med de viste fortegn,
      fx \(e=r-Hy\) for negativ sensorfeedback.
\item For hver blok skriver du output = blok gange input, fx \(u=Ce\) og \(y=Gu\).
      Indsæt derefter baglæns, indtil den efterspurgte variabel står i én ligning med
      inputtet.
\item Hvis diagrammet er standard negativ feedback, kan du kontrollere med
      \eqref{eq:standard-feedback}. Ellers må du ikke kopiere standardformlen; løs de
      skrevne ligninger algebraisk og faktoriser først til sidst.
\item Nævneren kommer fra de led der går rundt i feedbacksløjfen. Sidegrene uden for
      sløjfen hører til tælleren, medmindre de ændrer signalet inde i feedbackvejen.
\end{enumerate}

\paragraph{Formler}
For standard negativ feedback med forward branch \(C(s)G(s)\) og sensor \(H(s)\):
\begin{equation}
\label{eq:standard-feedback}
\frac{Y}{R}=\frac{CG}{1+CGH},\qquad
\frac{E}{R}=\frac{1}{1+CGH}.
\end{equation}

\paragraph{Symboler}
\(R,E,U,Y\) er reference, fejl, kontrolsignal og output; \(H\) er måleled.

\paragraph{Antagelser}
Formlerne gælder kun for det viste negative standardloop; ekstra feed-forward- eller
disturbancegrene lægges ind som særskilte bidrag.

\paragraph{Hurtigtjek}
Sæt en sidegren lig nul og kontroller at resultatet reduceres til et kendt enkelt loop.

\paragraph{Typiske fælder}
At sætte en forward sidegren ind i loopnævneren eller anvende unity-feedbackformlen på
en figur hvor gain ligger i målegrenen.

\paragraph{Python}
Ingen parser til diagrammet. Når du selv har reduceret til \(E/R=1/(1+KG)\), kan
\code{unity\_feedback\_step\_error} kontrollere slutværdien.

\subsection{Stationær referencefejl / final value, system type}

\paragraph{Brug når}
Der spørges om static/stationary/steady-state error efter step eller om systemtype.

\paragraph{Input fra opgaven}
Den udledte error transferfunktion, referenceform og bekræftelse af closed-loop-stabilitet.

\paragraph{Metode}
\begin{enumerate}
\item Udled først den rigtige error transferfunktion \(G_{er}(s)=E(s)/R(s)\) med
      blokdiagrammet i \S\ref{sec:feedback}. Brug ikke \(Y/R\), hvis opgaven spørger
      om fejl.
\item Indsæt referenceformen. For unit step er \(R(s)=1/s\); for et step med størrelse
      \(h_0\) er \(R(s)=h_0/s\). Dermed bliver \(E(s)=G_{er}(s)R(s)\).
\item Find polerne i den lukkede respons, dvs. nævneren i \(E(s)\) efter eventuelle
      forkortelser der er legitime i signalvejen. Final value theorem i
      \eqref{eq:steady-state-error} må kun bruges, når alle disse poler ligger strengt
      i \(\LHP\).
\item Beregn \(e_{\mathrm{ss}}\) med \eqref{eq:steady-state-error}: gang først med
      \(s\), forkort \(s\) fra step-inputtet, og sæt derefter \(s=0\).
\item Systemtypen bestemmes af open-loop transferfunktionen \(L(s)\): hvis \(L(s)\)
      har \(n\) rene integratorer, dvs. \(n\) faktorer \(s\) i nævneren ved origo,
      er systemtypen \(n\).
\end{enumerate}

\paragraph{Formler}
\begin{equation}
\label{eq:steady-state-error}
e_{\mathrm{ss}}=\lim_{s\to0}sE(s),\qquad
G_{er}(s)=S(s)=\frac1{1+L(s)}\quad\text{(negativ unity feedback)}.
\end{equation}
For stabil unity feedback giver type 0 endelig, normalt ikke-nul, fejl for step-input;
type 1 giver nul steady-state error for step og endelig fejl for ramp; type 2 giver
nul fejl for både step og ramp.
\begin{equation}
\label{eq:static-error-constants}
K_p^\ast=\lim_{s\to0}L(s),\qquad
K_v=\lim_{s\to0}sL(s),\qquad
K_a=\lim_{s\to0}s^2L(s).
\end{equation}
\begin{equation}
\label{eq:standard-input-errors}
e_{\mathrm{step}}=\frac1{1+K_p^\ast},\qquad
e_{\mathrm{ramp}}=\frac1{K_v},\qquad
e_{\mathrm{parabolic}}=\frac1{K_a}.
\end{equation}
Hvis konstanten er \(\infty\), er fejlen \(0\); hvis konstanten er \(0\), er fejlen
\(\infty\). Eksempel: \(G(s)=K/(s(s+3))\) er et type 1-system, fordi open-loop
transferfunktionen har én integrator, altså én faktor \(s\) i nævneren.

\paragraph{Symboler}
\(e_{\mathrm{ss}}\) er endelig fejl; systemtype er antal rene integratorer i open-loop
\(L(s)\).

\paragraph{Antagelser}
Final value theorem kræver asymptotisk stabilt closed loop for den betragtede respons.

\paragraph{Hurtigtjek}
F21 Q16: \(G(0)=2\), \(K_P=2\), derfor \(e_{\mathrm{ss}}=1/(1+4)=0.2\).

\paragraph{Typiske fælder}
At beregne \(Y/R\) når opgaven spørger om \(E/R\), eller at bruge en ikke-endelig DC-gain.

\paragraph{Python}
\code{unity\_feedback\_step\_error(num,den,proportional\_gain)} kontrollerer
standardloopet og afviser ustabil slutværdi; signalvejen er stadig manuel.
